%not a stand-alone document; intended for inclusion in a larger document

\section{Problem Geometries and Finite Element Meshes}\label{sec:geometry-and-mesh}

\textcolor{red}{\textbf{TODO}}: this should probably be restructured,
now that we actually have more than one geometry in question.

%-------------------------------------------------------------------------------
\subsection{Mesoporous Silica Membrane}\label{subsec:silica_membrane}

We study a silica membrane containing nanoscopic circular pores arranged in a two-dimensional lattice.
We desire to understand the rate of transport of different ions through the nanopores,
quantified in terms of an effective diffusion constant.
The pore geometry is variable.

We study both a body-centered rectangular lattice of pores,
as well as a face-centered lattice of pores.
The unit cell geometry has two planes of symmetry.
This symmetry is used to reduce the model to only one quarter of the unit cell.

The body-centered geometry is shown in Figure \ref{fig:body-intro}.
The face-centered geometry is shown in Figure \ref{fig:face-intro}.

\begin{figure}[H]
\centering
\includegraphics[width=1\textwidth]{body-intro.pdf}
\caption{Body-centered geometry}
\label{fig:body-intro}
\end{figure}

\begin{figure}[H]
\centering
\includegraphics[width=1\textwidth]{face-intro.pdf}
\caption{Face-centered geometry}
\label{fig:face-intro}
\end{figure}

The geometric variables are:

$\begin{array}{rcl}
Sx & = & 2 Lx =\text{Unit cell x-dimension} \\
Sy & = & 2 Ly =\text{Unit cell y-dimension} \\
Lx & = & \frac{Sx}{2} =\text{Model x-dimension} \\
Ly & = & \frac{Sy}{2} =\text{Model y-dimension} \\
R & = & \text{Pore radius} \\
tm & = & \text{Membrane thickness} = \text{Pore length} \\
H & = & \text{Distance from membrane surface to model boundary}
\end{array}$

\texttt{gmsh} is used to generate the finite element mesh.
The mesh is generated from solid geometry defined by its boundary surfaces.
These surfaces, in turn, are defined by their boundaries,
which are defined using points.

\textcolor{red}{\textbf{TODO}}: gmsh geometry (point numbers, etc)
